% Macros shared by the web (plasTeX/KaTeX) and print (pdf) builds.
% Morgan–Tian, *Ricci Flow and the Poincaré Conjecture* (arXiv:math/0607607).
%
% ── theorem-like environments ────────────────────────────────────────────────
% Numbered within chapters (1.1, 1.2, … / 2.1, …) to match the book. The hgraph
% sync parser keys a node off the environment NAME, so we use the parser's
% recognised names: theorem, proposition, lemma, corollary, definition, remark,
% example, claim, conjecture.
\newtheorem{theorem}{Theorem}[chapter]
\newtheorem{proposition}[theorem]{Proposition}
\newtheorem{lemma}[theorem]{Lemma}
\newtheorem{corollary}[theorem]{Corollary}
\newtheorem{claim}[theorem]{Claim}

\theoremstyle{definition}
\newtheorem{definition}[theorem]{Definition}
\newtheorem{example}[theorem]{Example}
\newtheorem{conjecture}[theorem]{Conjecture}
\newtheorem{remark}[theorem]{Remark}

% ── number sets (the book's \def\Ar{\mathbb R} etc.) ─────────────────────────
\newcommand{\Ar}{\mathbb{R}}
\providecommand{\Real}{\mathbb{R}}
\providecommand{\Complex}{\mathbb{C}}
\providecommand{\Zee}{\mathbb{Z}}
\providecommand{\ZZ}{\mathbb{Z}}
\providecommand{\RR}{\mathbb{R}}
\providecommand{\CC}{\mathbb{C}}
\providecommand{\QQ}{\mathbb{Q}}
\providecommand{\Q}{\mathbb{Q}}
\providecommand{\Cee}{\mathbb{C}}      % alternate spelling used once in Ch00
\providecommand{\HH}{\mathbb{H}}      % hyperbolic space H^n (book's \H clashes with amsmath)

% ── operators used across the two chapters ───────────────────────────────────
\DeclareMathOperator{\Ric}{Ric}
\DeclareMathOperator{\Rm}{Rm}
\DeclareMathOperator{\Hess}{Hess}
\DeclareMathOperator{\Vol}{Vol}
\DeclareMathOperator{\Tr}{Tr}
\DeclareMathOperator{\tr}{tr}
\DeclareMathOperator{\Jac}{Jac}
\DeclareMathOperator{\sn}{sn}
\DeclareMathOperator{\ct}{ct}
\DeclareMathOperator{\dist}{dist}
\DeclareMathOperator{\diam}{diam}
\DeclareMathOperator{\inj}{inj}
\DeclareMathOperator{\vol}{vol}
\DeclareMathOperator{\Sym}{Sym}
\DeclareMathOperator{\Hom}{Hom}
\DeclareMathOperator{\End}{End}
\DeclareMathOperator{\Spec}{Spec}
\providecommand{\exp}{\operatorname{exp}}
\DeclareMathOperator{\id}{id}

% ── plain-TeX primitives the book uses that KaTeX doesn't support directly ───
% \lefteqn hangs its argument to the left so a broken-across-lines equation
% aligns; KaTeX has no line-breaking model to hang against, so this just
% typesets the content in place (loses the hang, keeps the math readable).
\providecommand{\lefteqn}[1]{#1}

% ── convenience macros (book) ────────────────────────────────────────────────
\newcommand{\eps}{\varepsilon}
\newcommand{\To}{\longrightarrow}
\providecommand{\abs}[1]{\left\vert#1\right\vert}
\providecommand{\set}[1]{\left\{#1\right\}}
\providecommand{\seq}[1]{\left\langle#1\right\rangle}
\providecommand{\norm}[1]{\left\Vert#1\right\Vert}

% ── editorial commands neutralised for the blueprint ─────────────────────────
% (marginal notes and index entries carry no mathematical content here)
\providecommand{\note}[1]{}
\providecommand{\ii}[1]{\textit{#1}}

% ── cleveref fallback (plasTeX has no cleveref; pdflatex overrides this) ──────
\providecommand{\cref}[1]{\ref{#1}}
\providecommand{\Cref}[1]{\ref{#1}}
\providecommand{\crefrange}[2]{\ref{#1}--\ref{#2}}
\providecommand{\Crefrange}[2]{\ref{#1}--\ref{#2}}

% ── book accent redefinitions (chapters 3+) ──────────────────────────────────
% The source does \def\tilde{\widetilde} and \def\bar{\overline}; the wide forms
% are what the later chapters (L-length $\widetilde L$, reduced volume, etc.) rely
% on. Reproduce them so `\tilde{L}` / `\bar\tau` render as in the book.
\renewcommand{\tilde}{\widetilde}
\renewcommand{\bar}{\overline}
\providecommand{\part}{\partial}

% ── \mathcal shortcuts (the book's \def\cA..\cZ) ─────────────────────────────
\providecommand{\cA}{\mathcal{A}}\providecommand{\cB}{\mathcal{B}}
\providecommand{\cC}{\mathcal{C}}\providecommand{\cD}{\mathcal{D}}
\providecommand{\cE}{\mathcal{E}}\providecommand{\cF}{\mathcal{F}}
\providecommand{\cG}{\mathcal{G}}\providecommand{\cH}{\mathcal{H}}
\providecommand{\cI}{\mathcal{I}}\providecommand{\cJ}{\mathcal{J}}
\providecommand{\cK}{\mathcal{K}}\providecommand{\cL}{\mathcal{L}}
\providecommand{\cM}{\mathcal{M}}\providecommand{\cN}{\mathcal{N}}
\providecommand{\cO}{\mathcal{O}}\providecommand{\cP}{\mathcal{P}}
\providecommand{\cQ}{\mathcal{Q}}\providecommand{\cR}{\mathcal{R}}
\providecommand{\cS}{\mathcal{S}}\providecommand{\cT}{\mathcal{T}}
\providecommand{\cU}{\mathcal{U}}\providecommand{\cV}{\mathcal{V}}
\providecommand{\cW}{\mathcal{W}}\providecommand{\cX}{\mathcal{X}}
\providecommand{\cY}{\mathcal{Y}}\providecommand{\cZ}{\mathcal{Z}}

% ── a few more operators used from chapter 3 on ──────────────────────────────
\DeclareMathOperator{\RE}{Re}
\DeclareMathOperator{\IM}{Im}
\DeclareMathOperator{\Ker}{Ker}
\DeclareMathOperator{\Image}{Im}
\DeclareMathOperator{\supp}{supp}
\DeclareMathOperator{\Area}{Area}
\DeclareMathOperator{\Length}{Length}
\DeclareMathOperator{\Id}{Id}
\DeclareMathOperator{\scal}{scal}
